\subsubsection{Вычислительные системы}\label{subsec:ch1/sec2/sub4/sub1}

\fixme{
	пробить по источникам информации, добавить и плюс картинки для доклада
}

В автоматизированной системе управления НТТС могут использоваться следующие элементы вычислительных систем \cite{Tanenbaum, MunisterSoc, DigitalDesign} (рис.\ \cref{fig:sensors}):

\begin{itemize}[leftmargin=*]
	\item \textit{центральные процессоры (CPU)}:
	используются для выполнения общих задач, таких как управление операционной системой, обработка данных от датчиков и выполнение алгоритмов (процессоры Intel, AMD);

	\item \textit{графические процессоры (GPU)}:
	используются для задач машинного обучения и компьютерного зрения, таких как распознавание объектов, классификация изображений и обработка данных с камер (NVIDIA Drive AGX, NVIDIA Tesla, A100);

	\item \textit{тензорные процессоры (TPU)}:
	это специализированные процессоры для ускорения вычислений в нейронных сетях (Google TPU);

	\item \textit{системы на кристалле (SoC)}:
	более компактные решения, объединяющие CPU, GPU, TPU и другие компоненты на одном чипе (NVIDIA Xavier, NVIDIA Orin, Qualcomm Snapdragon Ride);

	\item \textit{программируемые логические интегральные схемы (FPGA)}:
	необходимы для задач, требующих высокой производительности и низкого энергопотребления, таких как обработка сигналов с лидаров и радаров (Xilinx Zynq, Intel FPGA);

	\item \textit{специализированные процессоры для автономных систем}:
	специальные для беспилотных транспортных средств, которые обеспечивают высокую производительность и энергоэффективность (NVIDIA Drive PX, Mobileye EyeQ);

	\item \textit{датчики и их интерфейсы}:
	лидары, радары, камеры, ультразвуковые датчики и GPS-модули подключаются к вычислительным системам через специализированные интерфейсы (CAN, Ethernet, PCIe);

	\item \textit{оперативная память}:
	высокоскоростная память для временного хранения данных, необходимых для обработки в реальном времени (DDR4, GDDR6);

	\item \textit{накопители данных}:
	используются для хранения больших объемов данных, таких как карты, модели машинного обучения и логи (NVMe SSD);

	\item \textit{коммуникационные модули}:
	обеспечение связи между компонентами системы, а также с внешними серверами (Ethernet, 5G-модули);

	\item \textit{микроконтроллеры}:
	для управления низкоуровневыми системами, такими как двигатели, тормоза и рулевое управление (ARM Cortex-M, STM32).
\end{itemize}

\nomenclature{CPU}{Central Processing Unit, центральный процессор\nomrefpage}
\nomenclature{GPU}{Graphics Processing Unit, графический процессор\nomrefpage}
\nomenclature{TPU}{Tensor Processing Unit, тензорный процессор\nomrefpage}
\nomenclature{SoC}{System on a Chip, электронная схема на одном кристалле\nomrefpage}
\nomenclature{FPGA}{field-programmable gate array, программируемая пользователем вентильная матрица\nomrefpage}

\begin{figure}[ht]
	\centering{
		\includegraphics[scale=0.6]{views/sensors.png}
	}
	\caption{Вариант схемы взаимодействия компонент системы автоматизированного управления в НТТС}\label{fig:sensors}
\end{figure}